\documentclass[11pt,a4paper]{article}
\usepackage[margin=22mm]{geometry}
\usepackage{amsmath,amssymb}
\pagestyle{empty}
\setlength{\parindent}{0pt}
\setlength{\parskip}{5pt}

\begin{document}
% References: tests/integrated/ion-energy-mms/manufactured.py and
% data-parallel_current{,_wall}/BOUT.inp.
\textbf{Three-dimensional manufactured-solution tests.}
We test the coupled ion total-energy, vorticity and electron parallel-momentum
equations on a Cartesian slab $0\leq X,Y,Z\leq L=10$, with periodic $Y$ and $Z$
and a uniform magnetic field parallel to $Y$.
All quantities below use Hermes normalised units, with reference density
$10^{18}\,\mathrm{m}^{-3}$, temperature $1\,\mathrm{eV}$ and magnetic field
$1\,\mathrm{T}$; lengths and time are normalised to the reference sound
gyroradius and inverse proton cyclotron frequency.
The two ion species have mass and charge numbers $(A_i,Z_i)=(2,1)$ and
$(A_j,Z_j)=(3,2)$; the electron mass is $1/1836$ in proton-mass units.
Define $k=\pi/L$, $a=kX$, $b=2kY$, $d=2kZ$, $\tau=t/5$,
$\theta=d-\tau$ and $W=\sin^2 a$. Both tests prescribe
\begin{align*}
 n_i &= 1+0.05\cos(2a)\cos(b-\tau)+0.03W\cos(d+\tau),\\
 n_j &= \tfrac14\bigl[1+0.05W\cos(2a)\sin\theta\cos(b+\tau)\bigr],\\
 v_i &= 0.10+0.02W\cos(2a)\sin(b-\tau)\cos\theta,\\
 v_j &= 0.07+0.015W\cos(2a)\cos(b+\tau)\sin\theta,\\
 T_e &= 1+0.05W\cos(2a)\sin\theta\cos(b-\tau),
 \qquad J_\parallel=0.03\cos(2a)\sin(b-\tau)\cos\theta.
\end{align*}
Here $n_s$ and $v_s$ denote species density and parallel velocity. Quasineutrality
and the prescribed current give $n_e=n_i+2n_j$ and
$v_e=(n_iv_i+2n_jv_j-J_\parallel)/n_e$.
The electrostatic and pressure potentials are
\begin{align*}
 A_\phi&=0.2[1+0.1\sin(b+\tau)],\qquad
 A_\psi=0.02[1+0.05\cos(b+\tau)],\\
 \phi&=A_\phi\sin a\cos\theta,\qquad
 \psi=1.2+0.04\cos(b-\tau)+A_\psi P(a)\sin\theta,\\
 P_0(a)&=[1-\cos(2a)]^2\quad\text{(zero wall curvature)},\qquad
 P_w(a)=1-\cos(2a)\quad\text{(finite wall curvature)}.
\end{align*}
Writing $\Phi=\phi+\psi$ and $\nabla_\perp=(\partial_X,\partial_Z)$, the remaining
manufactured fields follow from
\begin{align*}
 T_i=T_j&=\frac{2\psi}{2n_i+\tfrac32 n_j},\qquad
 p_i=n_iT_i,\quad p_j=n_jT_i,\quad p_e=n_eT_e,\\
 \Omega&=2\nabla_\perp^2\Phi,\qquad
 M_i=2n_iv_i,\quad M_j=3n_jv_j,\quad M_e=n_ev_e/1836,\\
 E_{\mathrm{ion}}&=\tfrac32(p_i+p_j)+n_iv_i^2+\tfrac32 n_jv_j^2
                  +|\nabla_\perp\Phi|^2.
\end{align*}
Here $M_s$ is parallel momentum density and $E_{\mathrm{ion}}$ includes ion
thermal, parallel kinetic and perpendicular kinetic energy. The analytic
polarisation potential, including parallel-current divergence, is
\[
 \chi=k^2 A_\phi A_\psi
       \sum_{n\in\{1,3,5\}}(c_n+d_n\sin^2\theta)\cos(na)
       -\frac{0.03}{8k}\cos(2a)\cos(b-\tau)\cos\theta,
\]
with coefficient vectors ordered as $n=1,3,5$:
\begin{align*}
 \boldsymbol c_0&=(16/17,-36/25,20/41),&
 \boldsymbol d_0&=(36/17,-78/25,42/41),\\
 \boldsymbol c_w&=(10/17,-18/25,0),&
 \boldsymbol d_w&=(14/17,-14/25,0).
\end{align*}
Continuum sources are obtained by substituting these fields into each equation
and subtracting its physical right-hand side from the exact time derivative.
Numerical potential damping is retained without cancellation by the sources.
At both radial walls, $\phi_{XX}=\psi_{XX}=0$ for $P_0$; for $P_w$,
$\phi_{XX}=0$ but $\psi_{XX}=4k^2A_\psi\sin\theta$.
This curvature concerns the potentials; magnetic curvature is zero in both
tests. The factor $W$ makes ion primitive wall values independent of $Z$,
avoiding a separate error from averaging wall transport weights.
We refine $(N_X,N_Y,N_Z)=(N,2N,N)$ and measure volume-weighted RMS ($L_2$) and
maximum ($L_\infty$) errors at $t=0.02$.
The zero-curvature test checks the full physical domain; the finite-curvature
test checks $X/L\in[1/4,3/4]$ and reports global errors to diagnose the Reynolds
wall contribution. Minimum accepted orders are $1.4$ for $L_2$ and $0.9$ for
$L_\infty$, consistent with MC slope limiting near extrema.
\end{document}
